\chapter*{本册导读：碳怎样活了起来}
\addcontentsline{toc}{chapter}{本册导读}
\markboth{本册导读}{}

第一册结束时，我们已经知道：原子怎样按电子结构排成周期表，原子之间怎样成键，反应的方向和速度由什么决定。那一套规则对食盐、铁钉、汽油和海水一视同仁。

第二册从一个特殊的元素开始：碳。它只有四个价电子，却因此能搭出链、环、支链和网格——整个有机化学，以及生命的全部分子材料，都建立在这只“四只手”的原子上。

这一册分三段旅程。第五篇先看碳怎样搭出分子世界：从石油到酒精，从阿司匹林到塑料，你会发现有机反应不需要一条条死背，官能团和电子的脾气会告诉你分子想干什么。第六篇进入生物化学：糖、脂质、蛋白质和酶怎样组成生命，ATP 怎样像货币一样流转能量，糖酵解、呼吸和光合作用怎样把食物、阳光和氧气接成一张网。第七篇把镜头再拉远一层：这些分子怎样组装成细胞——一座会维护自己、会听外界消息、会合作也会“叛变”的化学城市。

读这一册时，请常备第一册的两件工具：能量账（反应往能量更低、分散更多的方向走）和速率账（活化能决定快慢，催化剂只改路径不改总账）。你会发现，所谓“生命”，并不是这些账被违反了，而是被无比精巧地运用着。
